% GENERATED FILE. Edit canonical lessons, then run npm run generate:books.
% canonical-source-hash: fnv1a64:b39ed4fccfb198d1
% canonical-lessons: FR-C496-coudre, FR-C496-tricoter, FR-C496-pelote, FR-C496-broder, FR-C496-couture

\chapter{Fil et aiguille}
\label{ch:fr-fil-et-aiguille}
\hlchaptermodality{496}

\begin{chapteropening}
\textbf{By the end of this chapter:} \emph{I can talk about sewing, knitting and embroidery.}

Complete the last lesson of chapter 496: I can talk about sewing, knitting and embroidery.
\end{chapteropening}

\section[coudre]{coudre — to sew}

\label{lesson:FR-C496-coudre}



\begin{quote}
Before the new one: say the French for to betray, then the French for betrayal.
\end{quote}

\subsection*{You'll want to know: coudre}
\textbf{coudre} — \textquotedblleft{}to sew\textquotedblright{}.

\begin{quote}
Ma grand-mère m'a appris à coudre.
\end{quote}

\begin{quote}
\emph{My grandmother taught me to sew.}
\end{quote}

Irregular: \emph{je couds, nous cousons}; past \emph{j'ai cousu}. \emph{Coudre un bouton} — to sew on a button.

\begin{cousinweb}
\emph{coudre} is inherited from Vulgar Latin \emph{cōsere} (EtymDB), a form of Latin \emph{cōnsuere} \textquotedblleft{}to sew together\textquotedblright{}, built on \emph{suere} \textquotedblleft{}to sew\textquotedblright{} (Lewis \& Short).
\end{cousinweb}

\begin{grammarlens}[title={What you've built}]
One more word for telling, explaining and describing at B1.
\end{grammarlens}

\subsection*{Your turn}
\begin{itemize}
\raggedright
  \item \emph{Say it:} \emph{coudre}
  \item \emph{Say it:} \emph{coudre}, once more
  \item \emph{Say it:} a sentence of your own with \emph{coudre}
  \item \emph{Recall:} say the French for to betray, then the French for betrayal, then say \emph{coudre} again
\end{itemize}

\begin{tcolorbox}[breakable,colback=teal!4,colframe=teal!35!black,title={Before you move on}]
What does coudre mean? (\textquotedblleft{}To sew\textquotedblright{}.) Say it once more.
\end{tcolorbox}
\section[tricoter]{tricoter — to knit}

\label{lesson:FR-C496-tricoter}



\begin{quote}
Before the new one: say the French for betrayal, then the French for to sew.
\end{quote}

\subsection*{You'll want to know: tricoter}
\textbf{tricoter} — \textquotedblleft{}to knit\textquotedblright{}.

\begin{quote}
Mon père tricote une écharpe pour l'hiver.
\end{quote}

\begin{quote}
\emph{My father is knitting a scarf for the winter.}
\end{quote}

\emph{Le tricot} — knitting; also a knitted jumper.

\begin{cousinweb}
\emph{tricoter} goes back to Old French \emph{tricoter} (EtymDB), which once meant \textquotedblleft{}to beat with sticks\textquotedblright{}. Its origin is not settled; a Germanic source has been proposed.
\end{cousinweb}

\begin{grammarlens}[title={What you've built}]
One more word for telling, explaining and describing at B1.
\end{grammarlens}

\subsection*{Your turn}
\begin{itemize}
\raggedright
  \item \emph{Say it:} \emph{tricoter}
  \item \emph{Say it:} \emph{tricoter}, once more
  \item \emph{Say it:} a sentence of your own with \emph{tricoter}
  \item \emph{Recall:} say the French for betrayal, then the French for to sew, then say \emph{tricoter} again
\end{itemize}

\begin{tcolorbox}[breakable,colback=teal!4,colframe=teal!35!black,title={Before you move on}]
What does tricoter mean? (\textquotedblleft{}To knit\textquotedblright{}.) Say it once more.
\end{tcolorbox}
\section[la pelote]{la pelote — ball (of wool or string)}

\label{lesson:FR-C496-pelote}



\begin{quote}
Before the new one: say the French for to sew, then the French for to knit.
\end{quote}

\subsection*{You'll want to know: la pelote}
\textbf{la pelote} — \textquotedblleft{}ball (of wool or string)\textquotedblright{}.

\begin{quote}
Le chat joue avec une pelote de laine.
\end{quote}

\begin{quote}
\emph{The cat is playing with a ball of wool.}
\end{quote}

\emph{Une pelote de laine} — a ball of wool. \emph{La pelote basque} is a ball game.

\begin{cousinweb}
\emph{la pelote} is inherited from Vulgar Latin \emph{pilotta} \textquotedblleft{}a small ball\textquotedblright{} (EtymDB), built on Latin \emph{pila} \textquotedblleft{}a ball\textquotedblright{} (Lewis \& Short).
\end{cousinweb}

\begin{grammarlens}[title={What you've built}]
One more word for telling, explaining and describing at B1.
\end{grammarlens}

\subsection*{Your turn}
\begin{itemize}
\raggedright
  \item \emph{Say it:} \emph{la pelote}
  \item \emph{Say it:} \emph{la pelote}, once more
  \item \emph{Say it:} a sentence of your own with \emph{la pelote}
  \item \emph{Recall:} say the French for to sew, then the French for to knit, then say \emph{la pelote} again
\end{itemize}

\begin{tcolorbox}[breakable,colback=teal!4,colframe=teal!35!black,title={Before you move on}]
What does la pelote mean? (\textquotedblleft{}Ball (of wool or string)\textquotedblright{}.) Say it once more.
\end{tcolorbox}
\section[broder]{broder — to embroider}

\label{lesson:FR-C496-broder}



\begin{quote}
Before the new one: say the French for to knit, then the French for ball.
\end{quote}

\subsection*{You'll want to know: broder}
\textbf{broder} — \textquotedblleft{}to embroider\textquotedblright{}.

\begin{quote}
Elle a brodé des fleurs sur un coussin.
\end{quote}

\begin{quote}
\emph{She embroidered flowers on a cushion.}
\end{quote}

\emph{La broderie} — embroidery. Informally, to embellish a story.

\begin{cousinweb}
\emph{broder} goes back to Old French \emph{brosder}, a word of Germanic origin; its form was later drawn towards \emph{bord} \textquotedblleft{}edge\textquotedblright{}.
\end{cousinweb}

\begin{grammarlens}[title={What you've built}]
One more word for telling, explaining and describing at B1.
\end{grammarlens}

\subsection*{Your turn}
\begin{itemize}
\raggedright
  \item \emph{Say it:} \emph{broder}
  \item \emph{Say it:} \emph{broder}, once more
  \item \emph{Say it:} a sentence of your own with \emph{broder}
  \item \emph{Recall:} say the French for to knit, then the French for ball, then say \emph{broder} again
\end{itemize}

\begin{tcolorbox}[breakable,colback=teal!4,colframe=teal!35!black,title={Before you move on}]
What does broder mean? (\textquotedblleft{}To embroider\textquotedblright{}.) Say it once more.
\end{tcolorbox}
\section[la couture]{la couture — sewing; fashion design}

\label{lesson:FR-C496-couture}



\begin{quote}
Before the new one: say the French for ball, then the French for to embroider.
\end{quote}

\subsection*{You'll want to know: la couture}
\textbf{la couture} — \textquotedblleft{}sewing; fashion design\textquotedblright{}.

\begin{quote}
Je fais de la couture le samedi.
\end{quote}

\begin{quote}
\emph{I do sewing on Saturdays.}
\end{quote}

\emph{La haute couture} — high fashion. \emph{Une couture} is also a seam.

\begin{cousinweb}
\emph{la couture} is inherited, through Old French \emph{cousture}, from Vulgar Latin \emph{cōnsūtūra} \textquotedblleft{}a seam\textquotedblright{} (EtymDB), built on \emph{cōnsuere} \textquotedblleft{}to sew together\textquotedblright{}, from \emph{suere} (Lewis \& Short).
\end{cousinweb}

\begin{grammarlens}[title={What you've built}]
One more word for telling, explaining and describing at B1.
\end{grammarlens}

\subsection*{Your turn}
\begin{itemize}
\raggedright
  \item \emph{Say it:} \emph{la couture}
  \item \emph{Say it:} \emph{la couture}, once more
  \item \emph{Say it:} a sentence of your own with \emph{la couture}
  \item \emph{Recall:} say the French for ball, then the French for to embroider, then say \emph{la couture} again
\end{itemize}

\begin{tcolorbox}[breakable,colback=teal!4,colframe=teal!35!black,title={Before you move on}]
What does la couture mean? (\textquotedblleft{}Sewing; fashion design\textquotedblright{}.) Say it once more.
\end{tcolorbox}
